% Part: normal-modal-logic
% Chapter: tableaux
% Section: introduction

\documentclass[../../../include/open-logic-section]{subfiles}

\begin{document}

\olfileid{nml}{tab}{int}

\olsection{پېژندنه}

!!^{tableau}s د !!{signed
  formula}s ځانګړې (ښکته څانګې جوړوونکې)
ونې دي؛ يعنې داسې جوړې چې د رښتياوالي نښه ($\True$ يا
$\False$) او !!a{sentence} لري:
\[
\sFmla{\True}{!A} \text{ يا } \sFmla{\False}{!A}.
\]
!!^a{tableau} د يو شمېر \emph{فرضونو} څخه پيلېږي۔ هر بل
!!{signed formula} د يوې استنتاجي قاعدې په کارولو جوړېږي۔ ځينې
استنتاجي قاعدې د ونې يوې څوکې ته يو يا څو !!{signed formula}s
زياتوي؛ نورې دوه نوې څوکې جوړوي، چې دوه څانګې ترې پيدا کېږي۔
د قاعدو په پايله کښې داسې !!{signed formula}s راځي چې د هغو
!!{formula} د هغې !!{signed formula} له فارمول څخه لږ پېچلے وي چې قاعده
پرې پلې شوې وه۔ کله چې په يوه څانګه کښې دواړه
$\sFmla{\True}{!A}$ او $\sFmla{\False}{!A}$ وي، څانګه
\emph{تړلې} بولو۔ که د !!a{tableau} هره څانګه تړلې وي، ټول
!!{tableau} تړلے دے۔ تړلے !!{tableau} داسې !!a{derivation}
دے چې ښيي د !!{tableau} د پيل !!{signed formula}s سټ نۀ شي پوره
کېدے۔ په دې سره د~$\Proves$ اړيکه تعريفولے شو:
$\Gamma \Proves !A$ هله او يوازې هله چې داسې متناهي
سټ~$\Gamma_0 = \{!B_1, \dots, !B_n\} \subseteq \Gamma$ وي
چې د لاندې فرضونو لپاره تړلے !!{tableau} موجود وي:
\[
\{\sFmla{\False}{!A}, \sFmla{\True}{!B_1}, \dots, \sFmla{\True}{!B_n}\}.
\]

د مودالي منطقونو لپاره بايد د !!{signed
formula} مفهوم وغځوو
او داسې قاعدې هم زياتې کړو چې
\iftag{prvBox}{$\Box$\iftag{prvDiamond}{ او
    $\Diamond$}}{$\Diamond$} راونغاړي۔ په مودالي !!{tableau}s
کښې !!{formula}s د نښې ($\True$ يا $\False$) تر څنګ
\emph{مخکښونه}~$\sigma$ هم لري۔ مخکښونه د مثبتو صحيح عددونو
ناخالي لړۍ دي، يعنې $\sigma \in (\PosInt)^* \setminus
\{\emptyseq\}$۔ موږ دا مخکښونه له شاوخوا
$\tuple{\ }$ پرته ليکو، او جلا !!{element}s د~$,$ پر ځای
په~$.$ بېلوو۔ که $\sigma$ مخکښ وي، نو $\sigma.n$ له
$\sigma \concat \tuple{n}$ سره برابر دے؛ د بېلګې په توګه،
که $\sigma = 1.2.1$ وي، نو $\sigma.3$ برابر $1.2.1.3$
دے۔ نو د بېلګې په توګه،
\[
\sFmla{\True}{\Box !A \lif !A}[1.2]
\]
يو \emph{مخکښ‌لرونکے !!{signed formula}} دے (يا په لنډ ډول
\emph{مخکښ‌لرونکے !!{formula}})۔

د مفهوم له مخې، مخکښ په داسې مدل کښې د يوې نړۍ نوم دے چې ښايي
د !!a{tableau} د يوې څانګې !!{formula}s پوره کړي؛ که $\sigma$
د کومې نړۍ نوم وي، نو $\sigma.n$ د هغې نړۍ نوم دے چې له
($\sigma$ نومې) نړۍ څخه د لاسرسي وړ وي۔

\end{document}
